\documentclass[preview=true]{standalone}
\usepackage[siunitx, american]{circuitikz}
\usepackage[utf8]{inputenc}
\usepackage{newtxtext,newtxmath}
\usepackage[T1]{fontenc}

\begin{document}
% Define size/space
\def\loopsize{1cm}
\def\loopoffset{0.1cm}
% Define the loops
\def\myloops#1#2{
\begin{scope}[shift={#1}, scale=#2]
    % Draw the baseline
    \draw[line width=0.25mm] (-1,0) -- (1,0);
        % Draw the loops
    \draw[line width=0.25mm] (-\loopoffset,0) node [draw, circle, anchor=south, minimum size=\loopsize] (id) {};
    \draw[line width=0.25mm] (0,0) node [draw, circle, anchor=south, minimum size=\loopsize] (id) {};
    \draw[line width=0.25mm] (\loopoffset,0) node [draw, circle, anchor=south, minimum size=\loopsize] (id) {};
\end{scope}
}
\begin{circuitikz}
    \draw
      (0,0) node[left]{Port 1} [short] to [C] (2,0)
      (2,0) to [L, *-] (2,-2)
      (2,2) node[ground,rotate=180]{} to  [empty laser diode] (2,0) 
      (2,-2) [american current source, l^=$I_L$] to (2,-4) node[below]{$-V_L$};
      \myloops{(4,1)}{1};
    \draw
        (6,0) to [empty photodiode] (6,2) node[above]{$V_P$}
        (6,0) to [L,*-] (6,-2) node[ground]{}
        (6,0) to [C] (8,0) node[right]{Port 2}
    ; 
\end{circuitikz}
\end{document}